% ---------- status banner + abstract ----------
\begin{center}
{\small\textbf{RESEARCH PROTOCOL --- STUDY IN PROGRESS}}\\[2pt]
{\small Version 1.0.1, frozen 2026-09-14. No experimental results are reported in this document.}\\
{\small Repository: \url{https://github.com/iamcipherdev/VibeGuardBench}}
\end{center}

\vspace{4pt}
\begin{abstract}
\noindent
Prompt-driven AI application builders---systems marketed for ``vibe coding'' of full-stack products---can produce convincing applications from a single specification. What happens next is less clear: when deterministic engineering feedback identifies concrete defects in such an application, can the same builder repair those defects without breaking previously passing functionality? This protocol defines an empirical study of that question. Three frozen full-stack application specifications (appointment booking, multi-role dashboard, ordering and inventory) are each generated three times by two commercial AI application builders, yielding 18 initial applications. A deterministic oracle suite---clean install and build, static type checking and linting, browser-level functional tests, API-driven authorization probes, static security analysis, secret and dependency scanning, axe-core accessibility checks, and restart-based persistence probes---measures each application. For every failing check, a machine-generated, fixed-template feedback message is delivered to the builder, after which the \emph{entire} oracle suite re-executes, for up to three repair iterations. The study measures the repair success rate, the regression rate among previously passing checks, net quality change, and their trajectories across iterations and failure categories. The design, specifications, prompts, oracles, and analysis plan are frozen before data collection. This document reports no results: the experimental phase had not begun at the time of writing, and the only executed activity is a harness-validation pilot whose data is excluded from all research claims.
\end{abstract}

\section{Introduction}

The production of software by describing it in natural language has moved from
research demos to commercial products. Prompt-driven builders such as Lovable, Bolt,
and Replit generate complete web applications---frontend, backend, and
persistence---from a written specification, and the practice of assembling software
this way is widely known as vibe coding. Surveys of the area describe a fast-growing
tool landscape and identify reliability and security of the generated products as
open empirical problems \cite{ge2025vibe}.

Most evaluation effort concentrates on the first act: what does a builder produce
from a specification? Benchmarks such as WebGen-Bench score initially generated
websites with automated tests \cite{lu2025webgen}, FullStackBench scores generated
solutions across full-stack task domains \cite{byte2025fullstack}, and AppWorld
verifies agent programs against a simulated world of applications
\cite{trivedi2024appworld}. This emphasis leaves a second, practically important act
unmeasured: what happens when the generated application fails objective quality
checks and the builder is asked to fix it. In classical automated program repair,
the answer is sobering---patches that satisfy a test oracle frequently break behavior
elsewhere, a phenomenon documented as test overfitting \cite{smith2015cure},
patch overfitting \cite{petke2024patch}, and preservation violations
\cite{ismayilzada2023poracle}. Whether modern AI application builders exhibit
analogous behavior when repairing their own output, at the level of a whole
application with functional, security, accessibility, build, and persistence
concerns, has not, to our knowledge, been measured.

This protocol specifies such a measurement. The unit of study is an application
instance: generated by a commercial builder, frozen, measured by a deterministic
oracle suite, then repaired by the same builder under machine-generated feedback and
measured again, up to three times. The design makes three methodological commitments.
First, \emph{deterministic measurement}: every check produces the same outcome given
the same artifact; no LLM judge scores anything. Second, the \emph{whole-suite rule}:
after every repair, all checks re-execute, because regressions are only visible to
checks that run. Third, \emph{pre-freezing}: specifications, prompts, oracles,
metric definitions, and the analysis plan were committed publicly before any
experimental data exists.

Two boundaries of scope deserve emphasis at the outset. The study examines two
commercial builders and three bounded application tasks; its conclusions concern
those systems and tasks, not all AI coding tools. And because the builders' internal
models are not disclosed, the study compares \emph{systems}, not models.

\subsection{Contributions}

This document contributes (i) a frozen, preregistration-style research protocol for
measuring repair-and-regression dynamics of AI application builders; (ii) VibeGuardBench,
an open benchmark implementation---three frozen application specifications with
machine-checkable interface contracts, a 39-check deterministic oracle suite across
eight categories, an experiment runner, a feedback generator, and a metrics and
analysis pipeline; (iii) a related-work synthesis showing which parts of the question
adjacent benchmarks leave unanswered; and (iv) a harness-validation pilot that
demonstrates the full loop---generation, oracle execution, feedback, repair, and
re-measurement---on real artifacts, with its data explicitly excluded from research
claims.

\section{Background and Related Work}

\subsection{Benchmarks for initial generation}

WebGen-Bench evaluates LLM agents that generate websites from instructions, scoring
them with automated functional tests \cite{lu2025webgen}. A related benchmark,
FullStackBench, covers about
4{,}400 tasks across full-stack domains and scores solutions in a sandbox
\cite{byte2025fullstack}. AppWorld defines a controllable world of simulated apps
with state-based verification of agent behavior \cite{trivedi2024appworld}. DevBench
and its successor evaluate repository-scale generation across the development
lifecycle \cite{li2024devbench,kumarappan2026devbench}, and Commit0 measures
from-scratch library generation against given unit tests \cite{zhao2024commit0}.
These systems establish that automated, test-based evaluation of generated software
is feasible. They stop at initial generation: none of them feeds deterministic
failure information back to the generating system and measures what breaks.

\subsection{Repair against tests}

SWE-bench established the evaluation of issue-resolving patches with repository test
suites, separating FAIL\_TO\_PASS from PASS\_TO\_PASS outcomes \cite{jimenez2024swebench}.
A family of successors measures agentic repair \cite{yang2024sweagent}, simplified
non-agent pipelines \cite{xia2024agentless}, structure-aware fault localization
\cite{zhang2024autocoderover}, and the validity of the benchmark itself
\cite{aleithan2024swebenchplus}. Notably, Wang et al.\ find that patches reported as
solved can be incorrect beyond what the tests observe \cite{wang2026solved}. This
line evaluates repair of human-authored repositories against unit tests; the
artifact, the oracle, and the repair driver all differ from the setting studied here,
in which a commercial builder repairs its own generated application against a
multi-category QA suite.

\subsection{Regression and overfitting in repair}

That repaired software can be worse than its tests suggest is an old measurement.
Smith et al.\ show that most automatically produced patches are plausible rather than
semantically equivalent \cite{smith2015cure}. Poracle demonstrates that patches
violate preservation conditions beyond the failing test
\cite{ismayilzada2023poracle}; Petke et al.\ quantify the magnitude of patch
overfitting \cite{petke2024patch}; and Lou et al.\ show across two million candidate
patches that regression suites reject a substantial share of APR output
\cite{lou2024apr}. These studies concern classical search-based and semantic APR on
Java and C. The present study measures the same phenomenon---repair that breaks
passing behavior---in a modern setting: LLM-driven, builder-mediated repair of
full-stack applications, with regressions defined across functional, security,
accessibility, build, and persistence oracles rather than a unit-test suite alone.

\subsection{Iterative modification and self-correction}

CanItEdit evaluates instructional code editing, including successive edits where each
must preserve prior behavior, and finds editing capability degrades with instruction
complexity \cite{cassano2023canitedit}. Self-Refine improves outputs with
model-generated feedback \cite{madaan2023selfrefine}, while Huang et al.\ show that
intrinsic self-correction without external signals can degrade reasoning performance
\cite{huang2023selfcorrect}, and Reflexion uses environment signals across episodes
\cite{shinn2023reflexion}. Together this work suggests that the source of feedback
matters: deterministic, tool-generated feedback, as fixed in our design, removes the
self-judging confound that motivates the negative results.

\subsection{Security of AI-generated code}

Pearce et al.\ find that roughly 40\% of Copilot's completions in security-relevant
scenarios contain vulnerabilities \cite{pearce2022asleep}, and that LLMs can zero-shot
repair a non-trivial share of vulnerabilities in C/C++ snippets
\cite{pearce2023zeroshot}. Perry et al.\ show that developers assisted by an LLM write
less secure code and are more confident in it \cite{perry2023insecure}. Security
measurement of generated \emph{applications} in a repair loop, where the question is
whether security failures get repaired and at what regression cost, remains open; our
design includes static security analysis, secret detection, dependency auditing, and
behavioral authorization probes per iteration, with scanner findings manually
verified rather than taken at face value.

\subsection{Gap}

Our search of the venues and indexers listed in the repository's search log found no
empirical study combining commercial AI application builders, exported full-stack
applications, a deterministic multi-category QA oracle suite, machine-generated
repair feedback delivered back to the same builder, and whole-suite regression
measurement across iterations. The closest works are WebGen-Bench (initial
generation with tests, no repair loop) \cite{lu2025webgen}, the SWE-bench family
(unit-test oracles on human repositories) \cite{jimenez2024swebench}, and the APR
regression literature (pre-LLM tools) \cite{lou2024apr,ismayilzada2023poracle}.
Commercial comparisons of app builders exist in quantity as vendor and blog content,
but not as reproducible empirical studies. We therefore claim a \emph{coverage} gap,
not a first: the study's contribution is a careful measurement of a previously
unmeasured combination, with primary weight on measurement validity.

\section{Research Questions}

\begin{description}[leftmargin=1.6em]
\item[RQ1.] What proportion of objectively detected failures in AI-generated
full-stack web applications become passing after the same builder receives
deterministic, machine-generated feedback?
\item[RQ2.] How often do repair attempts introduce regressions, i.e., cause
previously passing checks to fail, measured per failure category?
\item[RQ3.] Do repair success and regression rates differ across check categories
(functionality, authentication/authorization, security, accessibility, build
reliability, database persistence, type correctness, maintainability)?
\item[RQ4.] How do repair success and regression behavior change over up to three
repeated repair cycles?
\item[RQ5 (exploratory).] Do the two studied builders differ in repairability and
regression profiles? The planned sample supports interval-estimated description
only.
\end{description}

RQ5 is deliberately demoted from confirmatory status: with 18 applications, a
powered comparison between builders is not attainable, and reporting one as if it
were would misstate the evidence.
